\subsection{质量剖面与最优配置的稳定性}
求得一个最优质量值之后，还应检查其附近损失曲线是否陡峭。图~\ref{fig:q3-quality-profiles} 在每一个固定质量上重新优化 $N,D$，绘制相对于联合候选最优值的可约损失增幅：
\begin{equation}
P_C(Q)=100\,\frac{F(Q;C)-L^*(C)}{L^*(C)-E}.
\label{eq:q3-quality-profile-regret}
\end{equation}
使用 $L^*-E$ 作为分母，是为了避免不可约常数 $E$ 掩盖可优化部分的差别。曲线上的非负增幅是确定性条件损失，并不是置信界。固定 $Q$ 后重新配置其余变量，也区别于只改变质量而将参数量和数据量锁死的单因素扰动。

\begin{figure}[!htbp]
\centering
\includegraphics[width=0.98\textwidth]{fig_q3_quality_profiles.pdf}
\caption{三档预算的质量剖面；每个质量点均重新优化参数量与数据量}
\label{fig:q3-quality-profiles}
\end{figure}

低预算下，候选最优质量只略高于基线；提高到更高质量将挤出大量训练数据，使总损失上升。中预算下最优值移向高质量端，但尚未完全触及上界；高预算下则以 $Q=1$ 为候选最优。剖面反映了上述变化的连续过程：最优点位置给出条件决策，而最优点附近的曲率说明轻微偏离的代价。若某一段较平坦，保留更多质量小数位并不意味着工程决策更加精确，还需考虑评分噪声和实施成本。

\subsection{总收益中规模配比与质量投入的分离}
将等数值基线、固定质量最优基线和联合优化值分别记为 $L_{\mathrm{eq}}$、$L_{Q_0}$ 与 $L^*$，则总降损存在精确的分解
\begin{equation}
L_{\mathrm{eq}}-L^*
=\underbrace{L_{\mathrm{eq}}-L_{Q_0}}_{\text{规模与数据配比调整}}
+\underbrace{L_{Q_0}-L^*}_{\text{放开质量决策的净改善}}.
\label{eq:q3-gain-decomp}
\end{equation}
两项都在同一个损失函数、同一预算和同一上下文条件下比较。第二项允许质量变化后重新优化 $N,D$，表示引入质量决策自由度的总净效益，并不等于只改变 $Q$ 的偏导收益。由于联合可行集包含 $Q=Q_0$ 的解，精确求解时该净改善应非负；数值实现则用容差处理机器精度量级的差异。

\begin{figure}[!htbp]
\centering
\includegraphics[width=0.88\textwidth]{fig_q3_gain_decomposition.pdf}
\caption{连续预算下的收益分解及质量净改善占固定质量可约损失的比例}
\label{fig:q3-gain-decomp}
\end{figure}

图~\ref{fig:q3-gain-decomp} 显示，相对朴素等数值方案的优势，不能全部归到数据质量上。低预算档放开质量决策仅带来约 $0.00018$ 的降损，大部分收益来自重新配置模型规模与 token 数。中预算的质量净改善扩大到约 $0.01812$，高预算为约 $0.01215$；绝对改变量并非随预算一直增加。由于高预算下总可约损失也下降，相同绝对量与相对量可以呈现不同变化趋势，应分别报告。

还需区分质量水平和质量成本份额。预算增大时，模型可能选择更高的 $Q$，但随着 $N$ 增大，$kN$ 相对于有界质量区间内的 $h(Q)$ 更大，从而 $C_Q/C=h/(kN+h)$ 可以下降。因此“高预算质量份额较低”不代表“高预算不重视质量”。在本模型中，达到质量上界后无法进一步增加 $Q$，新增预算主要流向规模和数据，质量份额下降具有明确的约束解释。

这些结论可用于形成条件配置规则：先以固定质量解析方案作为可行起点；检查质量基线处的单边收益是否足以覆盖数据量收缩；若足以覆盖，再沿质量剖面寻找内部候选并比较上界；最后用统一的实际成本式回算预算。该顺序能够区分基线选择、单位错误和数值搜索失败，但不会消除对真实损失函数、数据质量成本与上下文成本的标定需求。
